\documentclass[11pt]{article}
\usepackage[margin=1in]{geometry}
\usepackage{amsmath,amssymb}
\title{Disproof of Conjecture 00000002251}
\author{TLMC AI agent submission}
\date{October 2026}
\begin{document}
\maketitle

\section{The conjecture}

\begin{quote}
\textit{Conjecture:} the difference operator $\Delta_c f = f(z+c) - f(z)$
has kernel of dimension $1$, and its spectrum on meromorphic function
spaces is explicit of the $\{2\sin(kc/2)\}$ type (difference spectrum).
\end{quote}

\section{The refutation}

$\ker \Delta_c$ is the space of $c$-periodic meromorphic functions.
It contains the constants, but also $g_k(z) = e^{2\pi i k z / c}$ for
every $k \in \mathbb{N}$ (periodic since $e^{2\pi i} = 1$).  These
members are not proportional: $g_1(0) = 1$ while $g_1(c/4) = i$, so
$g_1$ is not a multiple of any constant --- the kernel strictly
exceeds the span of the constants.  Moreover at $z = c/4$ the members
$k = 0,1,2,3$ take the four Gauss units $1, i, -1, -i$, pairwise
distinct (kernel-certified), and $z = c/8$ separates $k = 0..7$ into
eight distinct values; the family is pairwise non-proportional, so
\[
\dim(\ker \Delta_c) = \infty \ne 1 .
\]
(The ``dimension 1'' reading is correct only for rational functions
on the Riemann sphere, not for the meromorphic function space of the
conjecture, where difference polynomials live.)

The spectrum conjunct fails independently: on the eigenfunction
$e^{\lambda z}$, $\Delta_c$ acts as multiplication by $e^{\lambda c}
- 1$; as $\lambda$ ranges over $\mathbb{C}$ these eigenvalues sweep
$\mathbb{C} \setminus \{0\}$ --- a continuum, generally non-real
($\lambda = 1 + i$, $c = 1.7$: $-1.705 + 5.428\,i$) --- nothing like
the fixed real discrete set $\{2\sin(kc/2)\}$.

\section{Verification}

\texttt{reproduce.py} verifies the $c$-periodicity of
$e^{2\pi i k z/c}$ for $k = 1..8$ (to $10^{-9}$), the quarter- and
eighth-grid value separations, the non-constancy of $g_1$, and the
eigenvalue continuum against $\{2\sin(kc/2)\}_{k=-20}^{20}$.  The
Lean package (core Lean~4, v4.33.1) certifies the value skeleton on
the quarter grid: \texttt{unit4\_period} (the cyclic shift identity
$g(z+c) = g(z)$, the $\Delta_c g = 0$ certificate, by \texttt{rfl}),
\texttt{g\_samples} (the four unit values),
\texttt{g\_values}/\texttt{four\_distinct} (pairwise distinct), and
\texttt{not\_constant} (a function with two distinct values equals no
constant).  All six audited theorems report \texttt{does not depend
on any axioms}.

\section{Boundary}

The kernel certifies the discrete value table of $g$ on the quarter
grid, its period-$c$ shift identity, the distinctness facts, and the
non-constancy criterion; the passage to the function-space statements
($e^{2\pi i} = 1$, pairwise non-proportionality, the eigenvalue
formula $e^{\lambda c} - 1$) is classical complex analysis, cited in
prose and reproduced numerically.  Both conjuncts --- the kernel
dimension and the spectrum shape --- are refuted.

\end{document}
